\textbf{Phần 1. Trắc nghiệm lựa chọn}
\setcounterpageref{ex}{0}

\begin{ex} %[5X11N64]
Biểu đồ hình quạt tròn thường dùng để biểu diễn
\choice
{Sự thay đổi số liệu theo thời gian}
{\True Tỉ số phần trăm của các phần trong toàn bộ tổng thể}
{Độ dài của các đoạn thẳng}
{Tọa độ các điểm trong không gian}
\loigiai{
Biểu đồ hình quạt tròn biểu diễn tỉ số phần trăm của các thành phần trong toàn bộ tổng thể $100\%$.\\
\textbf{Đáp án của câu hỏi là B}
}
\end{ex}

\begin{ex} %[5X11N64]
Tổng tỉ số phần trăm của tất cả các hình quạt trên một biểu đồ hình quạt tròn luôn bằng
\choice
{\True $100\%$}
{$50\%$}
{$180\%$}
{$360\%$}
\loigiai{
Toàn bộ hình tròn biểu diễn $100\%$.\\
\textbf{Đáp án của câu hỏi là A}
}
\end{ex}

\begin{ex} %[5X11H64]
Một biểu đồ hình quạt tròn biểu diễn thể loại phim yêu thích của học sinh gồm: Phim hoạt hình $50\%$, Phim khoa học $25\%$, Phim phiêu lưu $15\%$, Phim hài $10\%$. Thể loại phim chiếm một nửa hình tròn là
\choice
{Phim khoa học}
{Phim phiêu lưu}
{\True Phim hoạt hình}
{Phim hài}
\loigiai{
$50\%$ ứng với một nửa hình tròn, đó là phim hoạt hình.\\
\textbf{Đáp án của câu hỏi là C}
}
\end{ex}

\begin{ex} %[5X11H64]
Biểu đồ hình quạt tròn thể hiện tỉ lệ tham gia các môn thể thao của $200$ học sinh: Bóng đá $40\%$, Cầu lông $30\%$, Bơi lội $20\%$, Cờ vua $10\%$. Số học sinh tham gia môn Bóng đá là
\choice
{$60$ bạn}
{$70$ bạn}
{$90$ bạn}
{\True $80$ bạn}
\loigiai{
Số học sinh chơi bóng đá: $200 \times 40 : 100 = 80$ (bạn).\\
\textbf{Đáp án của câu hỏi là D}
}
\end{ex}

\begin{ex} %[5X11H64]
Biểu đồ quạt tròn biểu diễn cơ cấu cây trồng trong một khu vườn: Xoài $35\%$, Nhãn $25\%$, Cam $20\%$, Bưởi chưa rõ. Tỉ số phần trăm diện tích trồng Bưởi là
\choice
{$15\%$}
{\True $20\%$}
{$25\%$}
{$10\%$}
\loigiai{
Tỉ số phần trăm bưởi: $100\% - (35\% + 25\% + 20\%) = 20\%$.\\
\textbf{Đáp án của câu hỏi là B}
}
\end{ex}

\begin{ex} %[5X11H64]
Biểu đồ hình quạt tròn cho biết tỉ lệ học sinh khối 5 tham gia câu lạc bộ: Tiếng Anh $45\%$, Toán học $35\%$, Tin học $20\%$. Biết khối 5 có $160$ học sinh. Số học sinh tham gia CLB Tin học là
\choice
{$24$ bạn}
{$28$ bạn}
{\True $32$ bạn}
{$36$ bạn}
\loigiai{
$160 \times 20 : 100 = 32$ (bạn).\\
\textbf{Đáp án của câu hỏi là C}
}
\end{ex}

\begin{ex} %[5X11V64]
Khảo sát phương tiện đến trường của $300$ học sinh: Xe đạp chiếm $45\%$, Xe buýt chiếm $25\%$, Đi bộ chiếm $20\%$, Bố mẹ đưa đón chiếm $10\%$. Số học sinh đi xe đạp nhiều hơn số học sinh đi bộ là
\choice
{$60$ bạn}
{\True $75$ bạn}
{$80$ bạn}
{$90$ bạn}
\loigiai{
Chênh lệch tỉ lệ: $45\% - 20\% = 25\%$.\\
Số bạn chênh lệch: $300 \times 25\% = 75$ (bạn).\\
\textbf{Đáp án của câu hỏi là B}
}
\end{ex}

\begin{ex} %[5X11V64]
Một trang trại nuôi $500$ con gia cầm. Biểu đồ hình quạt cho thấy Gà chiếm $55\%$, Vịt chiếm $30\%$, Ngan chiếm $15\%$. Tổng số con Vịt và Ngan là
\choice
{\True $225$ con}
{$200$ con}
{$250$ con}
{$275$ con}
\loigiai{
Tỉ lệ Vịt và Ngan: $30\% + 15\% = 45\%$.\\
Số lượng: $500 \times 45\% = 225$ (con).\\
\textbf{Đáp án của câu hỏi là A}
}
\end{ex}

\begin{ex} %[5X11V64]
Biểu đồ quạt tròn cho biết kết quả xếp loại học tập của $120$ học sinh khối 5. Trong đó số học sinh Khá chiếm $40\%$, học sinh Giỏi chiếm $35\%$, còn lại là học sinh Trung bình. Số học sinh Trung bình là
\choice
{$24$ bạn}
{$36$ bạn}
{$42$ bạn}
{\True $30$ bạn}
\loigiai{
Tỉ lệ Trung bình: $100\% - (40\% + 35\%) = 25\%$.\\
Số học sinh Trung bình: $120 \times 25\% = 30$ (bạn).\\
\textbf{Đáp án của câu hỏi là D}
}
\end{ex}

\begin{ex} %[5X11V64]
Biểu đồ quạt biểu diễn chi tiêu trong tháng của gia đình bác An: Ăn uống $45\%$, Giáo dục $20\%$, Tiết kiệm $15\%$, Chi khác $20\%$. Biết số tiền dành cho Giáo dục là $4\,000\,000$ đồng. Tổng thu nhập hàng tháng của gia đình bác An là
\choice
{$16\,000\,000$ đồng}
{$24\,000\,000$ đồng}
{\True $20\,000\,000$ đồng}
{$25\,000\,000$ đồng}
\loigiai{
Tổng thu nhập: $4\,000\,000 : 20 \times 100 = 20\,000\,000$ (đồng).\\
\textbf{Đáp án của câu hỏi là C}
}
\end{ex}

\begin{ex} %[5X11V64]
Trên biểu đồ hình quạt tròn, một hình quạt có góc ở tâm là $90^\circ$ (chiếm $\dfrac{1}{4}$ hình tròn) thì biểu diễn tỉ lệ phần trăm là
\choice
{$20\%$}
{$30\%$}
{$50\%$}
{\True $25\%$}
\loigiai{
$\dfrac{90}{360} = \dfrac{1}{4} = 25\%$.\\
\textbf{Đáp án của câu hỏi là D}
}
\end{ex}

\begin{ex} %[5X11C64]
Biểu đồ hình quạt tròn cho biết cơ cấu bán hàng của siêu thị gồm $4$ nhóm: Thực phẩm ($40\%$), Đồ gia dụng ($25\%$), Thời trang ($20\%$), Hóa mỹ phẩm ($15\%$). Biết doanh thu từ Thực phẩm nhiều hơn Thời trang là $80$ triệu đồng. Tổng doanh thu của siêu thị là
\choice
{\True $400$ triệu đồng}
{$360$ triệu đồng}
{$450$ triệu đồng}
{$500$ triệu đồng}
\loigiai{
Hiệu tỉ lệ: $40\% - 20\% = 20\%$.\\
Tổng doanh thu: $80 : 20 \times 100 = 400$ (triệu đồng).\\
\textbf{Đáp án của câu hỏi là A}
}
\end{ex}

\textbf{Phần 2. Trắc nghiệm đúng sai}
\setcounterpageref{ex}{0}

\begin{ex} %[5X11H64]
Biểu đồ hình quạt biểu diễn cơ cấu các loại quả trong vườn: Cam $40\%$, Táo $30\%$, Bưởi $20\%$, Xoài $10\%$. Tổng số có $600$ cây. Xét tính đúng sai:
\choiceTF
{\True Số cây Cam trong vườn là $240$ cây.}
{\True Số cây Táo trong vườn là $180$ cây.}
{Số cây Bưởi ít hơn số cây Xoài là $60$ cây.}
{\True Tổng số cây Bưởi và Xoài là $180$ cây.}
\loigiai{
\begin{itemize}
    \item Cam: $600 \times 40\% = 240$ cây. Đúng.
    \item Táo: $600 \times 30\% = 180$ cây. Đúng.
    \item Bưởi ($20\% = 120$ cây) nhiều hơn Xoài ($10\% = 60$ cây). Khẳng định ít hơn là Sai.
    \item Bưởi + Xoài: $120 + 60 = 180$ cây. Đúng.
\end{itemize}
\begin{itemize}
    \item a) \textbf{Đúng}.
    \item b) \textbf{Đúng}.
    \item c) \textbf{Sai}.
    \item d) \textbf{Đúng}.
\end{itemize}
}
\end{ex}

\begin{ex} %[5X11V64]
Biểu đồ quạt tròn biểu diễn kết quả học lực của $240$ học sinh khối 5: Xuất sắc $30\%$, Tốt $45\%$, Hoàn thành $20\%$, Chưa hoàn thành $5\%$. Xét tính đúng sai:
\choiceTF
{\True Số học sinh Xuất sắc là $72$ bạn.}
{\True Số học sinh Tốt là $108$ bạn.}
{Số học sinh Chưa hoàn thành là $15$ bạn.}
{\True Số học sinh Chưa hoàn thành là $12$ bạn.}
\loigiai{
\begin{itemize}
    \item Xuất sắc: $240 \times 30\% = 72$ bạn. Đúng.
    \item Tốt: $240 \times 45\% = 108$ bạn. Đúng.
    \item Chưa hoàn thành: $240 \times 5\% = 12$ bạn $\neq 15$ bạn. Ý 3 Sai, Ý 4 Đúng.
\end{itemize}
\begin{itemize}
    \item a) \textbf{Đúng}.
    \item b) \textbf{Đúng}.
    \item c) \textbf{Sai}.
    \item d) \textbf{Đúng}.
\end{itemize}
}
\end{ex}

\textbf{Phần 3. Trả lời ngắn}
\setcounterpageref{ex}{0}

\begin{ex} %[5X11H64]
Một biểu đồ hình quạt tròn có 3 phần: Phần A chiếm $45\%$, Phần B chiếm $35\%$. Hỏi Phần C chiếm bao nhiêu phần trăm? (Chỉ ghi số).
\shortans{20}
\loigiai{
\textbf{Đáp số:} $20$\\
$100 - (45 + 35) = 20\%$.
}
\end{ex}

\begin{ex} %[5X11H64]
Biểu đồ quạt tròn cho biết có $25\%$ học sinh thích môn Bơi lội trên tổng số $160$ học sinh. Tính số học sinh thích môn Bơi lội.
\shortans{40}
\loigiai{
\textbf{Đáp số:} $40$\\
$160 \times 25 : 100 = 40$.
}
\end{ex}

\begin{ex} %[5X11V64]
Biểu đồ quạt tròn cho biết cơ cấu số sách thư viện: Sách giáo khoa $50\%$, Truyện tranh $30\%$, Sách tham khảo $20\%$. Biết thư viện có $1\,500$ cuốn sách. Tính số cuốn truyện tranh.
\shortans{450}
\loigiai{
\textbf{Đáp số:} $450$\\
$1\,500 \times 30 : 100 = 450$.
}
\end{ex}

\begin{ex} %[5X11V64]
Biểu đồ quạt cho biết $15\%$ số cây trong vườn là cây mít, tương ứng với $45$ cây mít. Tính tổng số cây trong vườn.
\shortans{300}
\loigiai{
\textbf{Đáp số:} $300$\\
$45 : 15 \times 100 = 300$.
}
\end{ex}

\begin{ex} %[5X11V64]
Biểu đồ quạt tròn biểu diễn tỉ lệ học sinh đi bộ đến trường là $35\%$, đi xe đạp là $45\%$. Tổng số học sinh toàn trường là $800$ em. Hỏi số học sinh đi xe đạp nhiều hơn đi bộ bao nhiêu em?
\shortans{80}
\loigiai{
\textbf{Đáp số:} $80$\\
$800 \times (45 - 35) : 100 = 80$.
}
\end{ex}

\begin{ex} %[5X11C64]
Một cửa hàng bán được $360\text{ kg}$ trái cây gồm Cam, Táo, Nho. Trên biểu đồ hình quạt tròn, phần diện tích biểu diễn Nho chiếm $\dfrac{1}{6}$ hình tròn, Cam chiếm $50\%$, còn lại là Táo. Tính số ki-lô-gam Táo đã bán.
\shortans{120}
\loigiai{
\textbf{Đáp số:} $120$\\
Nho: $360 : 6 = 60\text{ kg}$.\\
Cam: $360 \times 50\% = 180\text{ kg}$.\\
Táo: $360 - (60 + 180) = 120\text{ kg}$.
}
\end{ex}
